\section{Interlocking Rule Engine Module}
\label{sec:methodology:rule-engine}

The InterlockingRuleEngine is the central validation engine responsible for enforcing the rules that govern signal aspect transitions. It operates by loading declarative logic from structured JSON files, parsing control relationships and conditions independently of the codebase, and then evaluating requests against live infrastructure data obtained through the DatabaseManager. This separation of configuration from implementation provides significant flexibility, allowing safety rules to be updated or extended without recompilation and ensuring that domain experts can refine logic directly at the schema level.

When a signal aspect change is requested, the engine applies a multi-stage evaluation process. It determines whether the signal is independent or subject to control by other signals, validates the status of track circuits and point machines, and checks that all defined conditions are satisfied. The system then aggregates these results according to the defined control mode, applying either strict or permissive logic as specified. The outcome is returned as a validation result object that records the decision, the reasoning behind it, and any affected infrastructure, ensuring that operators and auditors retain full visibility into the decision-making process.

The design of the engine emphasizes modularity, with dedicated components for parsing, evaluation, and database access. Diagnostic logging is embedded throughout, allowing each validation step to be reconstructed for debugging or forensic analysis. By combining declarative rule definition with modular execution, the InterlockingRuleEngine ensures that interlocking logic remains transparent, adaptable, and verifiable—qualities that distinguish it from the rigid, hardwired behavior of relay or PLC-based systems.